\section*{Chapitre \ref{ch:ml:cross-sectional-deep-models} --- Modèles profonds en coupe transversale}

\begin{solution}{exo:ml:cross-sectional-deep-models:1}
$\Gamma$ est $11\times3$~: 33 paramètres, quel que soit le nombre d'actions. ACP statique~: $300\times3 = 900$ charges, et une action
sans historique n'en a aucune.
\end{solution}

\begin{solution}{exo:ml:cross-sectional-deep-models:2}
$(0.5\times2 + (-0.5)(-1) + 1.0\times3)/3 = 1.5~\%$. Une somme sur les actions ne dépend pas de leur ordre.
\end{solution}

\begin{solution}{exo:ml:cross-sectional-deep-models:3}
Encodage one-hot vers 32 unités~: $30\times32 = 960$ poids. Un plongement unidimensionnel~: 30 valeurs, plus 32 poids reliant le
plongement à la couche, soit 62 au total.
\end{solution}

\begin{solution}{exo:ml:cross-sectional-deep-models:4}
Les charges de l'ACP sont les moyennes des bêtas de chaque action sur les années d'apprentissage, alors que les bêtas évoluent avec les caractéristiques~:
aux mois de test, les bêtas courants des actions, et donc leurs rendements espérés, n'ont guère de rapport avec ces moyennes
(corrélation 0.03). Le R² total provient surtout du premier facteur, un facteur de marché dont les bêtas sont proches de un, que toute
charge capte.
\end{solution}

\begin{solution}{exo:ml:cross-sectional-deep-models:5}
Le R² prédictif par rapport à zéro dépend du niveau de la prévision relativement au rendement moyen réalisé~; sur soixante
mois, les moyennes réalisées des facteurs s'écartent des primes, et un modèle dont la prévision moyenne les épouse mieux
obtient un score supérieur à celui de la vérité (1.23~\% contre 1.09~\%). Rapporter le R² total, la corrélation des rendements espérés avec la
vérité (sur données simulées) ou avec les rendements réalisés (l'IC), et leur dispersion selon les périodes.
\end{solution}

\begin{solution}{exo:ml:cross-sectional-deep-models:6}
L'ajustement dans l'échantillon croît avec la flexibilité et le nombre de facteurs~; il faut comparer les deux modèles hors échantillon, à même nombre
de facteurs, et avec des graines aléatoires pour l'\omterm{def:ml:cross-sectional-deep-models:ae}{autoencodeur}. Comparer dix facteurs à trois reviendrait aussi à comparer des modèles différents.
\end{solution}

\begin{solution}{exo:ml:cross-sectional-deep-models:7}
\texttt{embedding(emb\_init=None)}~: la corrélation du plongement avec les primes vaut 0.05 et celle des prévisions avec
la vérité 0.53 (pas mieux que sans l'industrie). Avec un écart type initial de un, les valeurs initiales aléatoires du
plongement sont bien plus grandes que l'effet des primes sur la perte ne peut les déplacer avant que l'\omterm{def:ml:trees-and-boosting:early}{arrêt précoce} (early stopping) ne mette fin à l'apprentissage~; partant
de valeurs proches de zéro, ce qu'apprend le plongement est tout ce qu'il contient.
\end{solution}

\begin{solution}{exo:ml:cross-sectional-deep-models:8}
Avec $\Gamma = I$ et $K = L$, les facteurs du mois $t$ sont $(Z_t^\top Z_t)^{-1}Z_t^\top r_{t+1}$, soit les coefficients des moindres carrés
des rendements de ce mois sur ses caractéristiques~: première étape de Fama et MacBeth. L'IPCA avec $K < L$ contraint
ces coefficients à un sous-espace de dimension $K$ décrit par $\Gamma$.
\end{solution}

\begin{solution}{pb:ml:cross-sectional-deep-models:1}
\begin{enumerate}
\item Total 31.3~\%, prédictif 1.09~\%.
\item Les rendements spécifiques (8~\% par mois) représentent la plus grande part de la variance, et aucun modèle factoriel ne les explique.
\item 0.6~\%, 0.4~\% et 0.3~\% par mois~; ils font surtout dépendre le R² prédictif du niveau moyen.
\item La moitié de la variance en coupe transversale de chaque bêta provient de carrés, de produits et de valeurs absolues que nulle fonction linéaire
des caractéristiques ne reproduit.
\item ACP 22.0~\%, 22.5~\%, 23.0~\%~; IPCA 25.1~\%, 27.6~\%, 28.1~\%~; \omterm{def:ml:cross-sectional-deep-models:ae}{autoencodeur} 26.9~\%, 30.4~\%, 30.8~\%.
\item ACP 0.03 et 0.00~; IPCA 0.71 et $-0.03$~; \omterm{def:ml:cross-sectional-deep-models:ae}{autoencodeur} 0.95 et 0.59.
\item Peu~: 30.4--30.6~\% au total et corrélation de 0.95--0.96 sur trois graines aléatoires.
\item Ses bêtas sont linéaires en les caractéristiques, et la partie non linéaire est orthogonale aux fonctions linéaires de celles-ci par
construction.
\item Les bêtas sont une seule fonction des caractéristiques pour chaque action, si bien que partager les paramètres entre actions ne coûte rien.
\item R² de test de 0.004~\% à 0.15~\%, corrélation avec la vérité de 0.52 à 0.75~; le plongement est corrélé à 0.86 avec
les primes, contre 0.88 pour les moyennes brutes par secteur.
\item Son initialisation~: avec un écart type égal à un, il n'a jamais appris les primes.
\item En entrée du réseau de facteurs, les portefeuilles gérés $Z_t^\top r_{t+1}/n$.
\item \emph{Résultat nommé}~: à trois facteurs, R² total de 22.5~\% (ACP), 27.6~\% (IPCA) et 30.4~\% (\omterm{def:ml:cross-sectional-deep-models:ae}{autoencodeur}) contre la
vérité, 31.3~\%~; R² prédictif 1.01~\%, 1.11~\% et 1.23~\% (vérité 1.09~\%)~; l'\omterm{def:ml:cross-sectional-deep-models:ae}{autoencodeur} retrouve 0.59 de la non-linéarité
injectée en corrélation, IPCA et ACP n'en retrouvent aucune.
\item Les bêtas de l'\omterm{def:ml:cross-sectional-deep-models:ae}{autoencodeur} (ou ceux de l'IPCA lorsque la non-linéarité est faible)~: les expositions aux facteurs doivent être celles des
caractéristiques courantes.
\item Sur des mois mis de côté, selon le R² total, en s'arrêtant là où un facteur de plus n'aide plus.
\item Normaliser par rang, imputer les valeurs manquantes par date, et s'attendre à ce que l'avantage de l'\omterm{def:ml:cross-sectional-deep-models:ae}{autoencodeur} diminue avec le signal.
\item Rien à tirer de l'ACP~; les bêtas issus de ses caractéristiques viennent de l'IPCA et de l'\omterm{def:ml:cross-sectional-deep-models:ae}{autoencodeur}.
\item Graines aléatoires, réglage, durée d'apprentissage et interprétabilité~: le $\Gamma$ de l'IPCA peut se lire.
\item Comparer sur des mois mis de côté, à même nombre de facteurs, sur plusieurs graines aléatoires et plusieurs périodes.
\item Comment les expositions d'une action aux risques communs dépendent de ce qu'elle est.
\end{enumerate}
\end{solution}

\begin{solution}{iq:ml:cross-sectional-deep-models:1}
L'ACP attribue à chaque action des charges fixes estimées à partir de son historique de rendements~; l'IPCA fait des charges des fonctions de caractéristiques observables,
caractéristiques, si bien qu'elles varient avec elles, existent pour les nouvelles actions et comptent beaucoup moins de paramètres.
\iqlookfor{bêtas variables dans le temps, pilotés par les caractéristiques.}
\end{solution}

\begin{solution}{iq:ml:cross-sectional-deep-models:2}
Un réseau de bêtas associe les caractéristiques de chaque action à $K$ bêtas~; un réseau de facteurs associe les portefeuilles gérés selon les caractéristiques du
mois à $K$ facteurs~; leur produit reconstruit les rendements, et la perte est l'erreur de reconstruction.
\iqlookfor{les deux entrées et la raison pour laquelle l'entrée des facteurs est un ensemble de portefeuilles.}
\end{solution}

\begin{solution}{iq:ml:cross-sectional-deep-models:3}
Lorsque des actions présentent un comportement idiosyncratique persistant que les caractéristiques ne captent pas, et que chacune dispose d'un historique suffisant~; il peut
mémoriser le bruit, échoue pour les nouveaux identifiants, et fuit par les identifiants qui changent (réattribution de tickers, fusions). Initialiser le
réseau avec de petites valeurs, puis le régulariser.
\iqlookfor{le cas d'usage, les modes de défaillance et les précautions pratiques.}
\end{solution}

\begin{solution}{iq:ml:cross-sectional-deep-models:4}
Le R² total compare la qualité avec laquelle les bêtas expliquent les rendements étant donné les facteurs~; le R² prédictif sur une courte période mêle la
coupe transversale au niveau des primes et à la chance de la période. Utiliser le R² total et des mesures en coupe transversale (IC, écarts
de déciles) pour comparer.
\iqlookfor{ce que mesure chacun et son bruit.}
\end{solution}

\begin{solution}{iq:ml:cross-sectional-deep-models:5}
Agréger avec une fonction invariante par permutation~: moyennes de caractéristiques, rendements de portefeuilles pondérés par les caractéristiques,
ou un encodeur d'ensemble appris $\rho(\sum_i\phi(x_i))$.
\iqlookfor{invariance par permutation et résumés de taille fixe.}
\end{solution}

\begin{solution}{iq:ml:cross-sectional-deep-models:6}
Étant donné $\Gamma$, les facteurs de chaque mois sont la régression des rendements sur $Z_t\Gamma$~; étant donné les facteurs, $\Gamma$ intervient linéairement
à travers $z\otimes f$, et ses équations normales sont $\sum_t(Z_t^\top Z_t\otimes f_{t+1}f_{t+1}^\top)\operatorname{vec}(\Gamma^\top)
= \sum_tZ_t^\top r_{t+1}\otimes f_{t+1}$~; alterner, en normalisant $\Gamma^\top\Gamma = I$.
\iqlookfor{les deux étapes et la structure de Kronecker.}
\end{solution}
